\documentclass{article}
\usepackage{amsmath}
\usepackage{graphicx}
\usepackage{setspace}
\begin{document}
	\begin{itemize}
		\item [3.]
		Its dual is:
		\begin{equation*}
			\begin{aligned}
				\text{min}&&4&\pi_1&+&&3&\pi_2&&&&\\
				\text{s.t.}&&5&\pi_1&-&&2&\pi_2&\geq&&4&&\text{(Constraint 1)}\\
				&&-&\pi_1&+&&&\pi_2&\geq&&-1&&\text{(Constraint 2)}\\
				&&2&\pi_1&+&&&\pi_2&\geq&&3&&\text{(Constraint 3)}\\
				&&&\pi_1&&&&&\geq&&0&\\
				&&&&&&&\pi_2&\geq&&0&
			\end{aligned}
		\end{equation*}
		Suppose that $\pi=(\frac{4}{3},\frac{1}{3})$ is optimal. With this solution, $\pi_1\neq 0, \pi_2\neq 0,$ constraint 2 and 3 are active, while 1 is not.\\
		So in any optimal solution to the primal problem, $x_1=0$, and both the two constraints are active. Solve them and we get $x=(0,\frac{2}{3},\frac{7}{3})$, and in both primal and dual problems, the objective functions equal to $\frac{19}{3}$.\\
		So the given $\pi$ does be an optimal solution to the dual problem, and the $x$ obtained above is an optimal solution to the primal problem.
	\end{itemize}
\end{document}